% Part: proof-theory
% Chapter: natural-deduction
% Section: rules-N2

\documentclass[../../../include/open-logic-section]{subfiles}

\begin{document}

\begin{table}
\begin{tabular}{@{}c@{}c@{}}
\Axiom$x: !A, \Gamma \fCenter !B$
\RightLabel{\Intro{\lif}}
\UnaryInf$\Gamma \fCenter !A \lif !B $
\DisplayProof
&
\Axiom$\Gamma_1 \fCenter !A \lif !B$
\Axiom$\Gamma_2 \fCenter !A$
\RightLabel{\Elim{\lif}}
\BinaryInf$\Gamma_1, \Gamma_2 \fCenter !B$
\DisplayProof
\\[4ex]
\Axiom$\Gamma_1 \fCenter !A$
\Axiom$\Gamma_2 \fCenter !B$
\RightLabel{\Intro{\land}}
\BinaryInf$\Gamma_1, \Gamma_2 \fCenter !A \land !B $
\DisplayProof
&
\begin{tabular}[b]{@{}l@{}}
\Axiom$\Gamma \fCenter !A \land !B$
\RightLabel{\Elim{\land}[1]}
\UnaryInf$\Gamma \fCenter !A$
\DisplayProof
\\[3ex]
\Axiom$\Gamma \fCenter !A \land !B$
\RightLabel{\Elim{\land}[2]}
\UnaryInf$\Gamma \fCenter !B$
\DisplayProof\\[3ex]
\end{tabular}
\\[4ex]
\Axiom$\Gamma \fCenter !A$
\RightLabel{\Intro{\lor}[1]}
\UnaryInf$\Gamma \fCenter !A \lor !B$
\DisplayProof
&
\Axiom$ \Gamma \fCenter !B$
\RightLabel{\Intro{\lor}[2]}
\UnaryInf$\Gamma \fCenter !A \lor !B$
\DisplayProof\\[3ex]
\multicolumn{2}{c}{
\begin{tabular}[b]{@{}l@{}}
\Axiom$\Gamma_1 \fCenter !A \lor !B$
\Axiom$x:!A, \Gamma_2 \fCenter !C$
\Axiom$y:!B, \Gamma_3 \fCenter !C$
\RightLabel{\Elim{\lor}}
\TrinaryInf$\Gamma_1, \Gamma_2, \Gamma_3 \fCenter !C$
\DisplayProof
\end{tabular}
}
\\[4ex]
\Axiom$x:!A, \Gamma \fCenter \lfalse$
\RightLabel{\Intro{\lnot}}
\UnaryInf$ \Gamma \fCenter \lnot !A $
\DisplayProof
&
\Axiom$\Gamma_1 \fCenter \lnot !A $
\Axiom$\Gamma_2 \fCenter !A $
\RightLabel{\Elim{\lnot}}
\BinaryInf$\Gamma_1, \Gamma_2 \fCenter \lfalse$
\DisplayProof
\\[4ex]
\Axiom$\Gamma \fCenter !A(a) $
\RightLabel{\Intro{\lforall}}
\UnaryInf$\Gamma \fCenter \lforall[x][!A(x)]$
\DisplayProof
&
\Axiom$\Gamma \fCenter \lforall[x][!A(x)]$
\RightLabel{\Elim{\lforall}}
\UnaryInf$\Gamma \fCenter !A(t)$
\DisplayProof
\\[4ex]
\Axiom$\Gamma \fCenter !A(t) $
\RightLabel{\Intro{\lexists}}
\UnaryInf$\Gamma \fCenter \lexists[x][!A(x)]$
\DisplayProof
&
\Axiom$\Gamma_1 \fCenter  \lexists[x][!A(x)]$
\Axiom$x: !A(a), \Gamma_2 \fCenter !C$
\RightLabel{\Elim{\lexists}}
\BinaryInf$\Gamma_1, \Gamma_2 \fCenter !C$
\DisplayProof
\\[4ex]
\Axiom$\Gamma \fCenter \lfalse$
\RightLabel{\FalseInt}
\UnaryInf$\Gamma \fCenter !A$
\DisplayProof
&
\Axiom$x: \lnot !A, \Gamma \fCenter \lfalse$
\RightLabel{\FalseCl}
\UnaryInf$\Gamma \fCenter !A$
\DisplayProof
\end{tabular}
\caption{د \Log{N2c} او \Log{N2i} قاعدې۔ په $\Intro{\forall}$ او
$\Elim{\exists}$ کښې !!{constant}~$a$ بايد په پاييز سېکوېنټ کښې
نۀ راځي؛ د موجود کميت د ايستلو په لويه مقدمه کښې هم دغه ثابت
نۀ شي راتلے۔ $\Gamma_i$~د نښه لرونکو \psOblique{!!{formula}s}~$x :
!A$ سټونه دي۔ د \Intro{\lforall} او \Elim{\lexists} قاعدې بايد
د ځانګړي متغير شرط پوره کړي: $a$ بايد په پاييز سېکوېنټ کښې
نۀ راځي او د موجود کميت د ايستلو په لويه مقدمه کښې هم نۀ راځي۔
\Log{N2i} د \FalseCl لرونکے نۀ دے۔}
\textbf{سرچينه‌يي سمون OLCMP-139:} د موجود کميت په لويه
مقدمه کښې هم ځانګړے ثابت نۀ راځي؛ يوازې د پاييز سېکوېنټ
تازګي کافي نۀ ده۔ د N1 اصلي شرط دا اړين قيد هم لري۔
\ollabel{tab:N2}
\end{table}

\end{document}